% Part 4: Experimental Setup and Implementation

\section{Experimental Setup and Implementation}

\subsection{Simulation Environment}

We implemented a comprehensive LEO satellite network simulator in Python using the following components:

\begin{itemize}
\item \textbf{Orbital Mechanics}: Circular orbit model with Walker-Delta constellation configuration
\item \textbf{Network Stack}: Custom discrete-event simulator for packet forwarding and queue management
\item \textbf{DRL Framework}: Stable-Baselines3 implementation of PPO \cite{raffin2021stable}
\item \textbf{Prediction Models}: PyTorch implementation of LSTM and GRU networks
\end{itemize}

\subsection{Constellation Configuration}

Table \ref{tab:constellation} summarizes the satellite constellation parameters.

\begin{table}[h]
\centering
\caption{LEO Satellite Constellation Parameters}
\label{tab:constellation}
\begin{tabular}{|l|c|}
\hline
\textbf{Parameter} & \textbf{Value} \\
\hline
\hline
Total Satellites & 248 \\
Orbital Planes & 28 (3 layers) \\
Satellites per Plane & Variable (6-10) \\
Altitude - Layer 1 & 500 km \\
Altitude - Layer 2 & 800 km \\
Altitude - Layer 3 & 1200 km \\
Inclination & 53° / 60° / 70° \\
Orbital Period & 94-113 minutes \\
Max ISL Distance & 5000 km \\
Max Neighbors & 8 \\
ISL Capacity & 300 Mbps \\
Queue Capacity & 50 packets \\
\hline
\end{tabular}
\end{table}

The constellation provides global coverage with inter-plane and intra-plane ISLs. Figure \ref{fig:3d_constellation} shows the 3D visualization of the constellation.

\begin{figure}[h]
\centering
\includegraphics[width=0.48\textwidth]{outputs/3d_viz/constellation_global.png}
\caption{3D visualization of 248-satellite LEO constellation with three orbital layers. Different colors represent different altitude layers.}
\label{fig:3d_constellation}
\end{figure}

\subsection{Traffic Configuration}

Traffic is generated according to three slices with parameters shown in Table \ref{tab:traffic}.

\begin{table}[h]
\centering
\caption{Traffic Slice Configuration}
\label{tab:traffic}
\begin{tabular}{|l|c|c|c|}
\hline
\textbf{Parameter} & \textbf{URLLC} & \textbf{eMBB} & \textbf{mMTC} \\
\hline
\hline
Priority & 1 (High) & 2 (Med) & 3 (Low) \\
Weight & 35\% & 35\% & 30\% \\
Packet Size & 256 B & 1500 B & 128 B \\
Max Latency & 25 ms & 50 ms & 200 ms \\
Target PDR & 99.999\% & 99.9\% & 99\% \\
Bandwidth Req. & 10 Mbps & 100 Mbps & 1 Mbps \\
\hline
\end{tabular}
\end{table}

Traffic is generated at 50 packets/second globally with random source-destination pairs following a realistic geographic distribution.

\subsection{Prediction Model Configuration}

Both LSTM and GRU models share the following architecture:

\begin{itemize}
\item Input features: 7 (queue, utilization, bandwidth, drops, delay, arrival rate, success rate)
\item Hidden size: 128
\item Number of layers: 2
\item Dropout: 0.2
\item Lookback window: 50 timesteps
\item Prediction horizon: 10 timesteps
\item Batch size: 32
\item Learning rate: 0.001
\item Optimizer: Adam
\end{itemize}

Training data was collected from 5 episodes of 150 timesteps each, totaling 1,693,896 link records. Data was split 80\%-20\% for training and validation. Figure \ref{fig:training_data} shows the data collection process.

\begin{figure}[h]
\centering
\includegraphics[width=0.48\textwidth]{training_data_collection.png}
\caption{Traffic data collection showing queue length evolution across multiple links over time. The diverse patterns demonstrate the need for adaptive prediction.}
\label{fig:training_data}
\end{figure}

\subsection{PPO Hyperparameters}

Table \ref{tab:ppo_params} lists the PPO training hyperparameters.

\begin{table}[h]
\centering
\caption{PPO Hyperparameters}
\label{tab:ppo_params}
\begin{tabular}{|l|c|}
\hline
\textbf{Parameter} & \textbf{Value} \\
\hline
\hline
Learning Rate & $10^{-4}$ \\
Steps per Update & 8192 \\
Batch Size & 1024 \\
Number of Epochs & 10 \\
Discount Factor ($\gamma$) & 0.995 \\
GAE Lambda ($\lambda$) & 0.98 \\
Clip Range ($\epsilon$) & 0.15 \\
Value Coefficient ($c_1$) & 0.7 \\
Entropy Coefficient ($c_2$) & 0.005 \\
Max Gradient Norm & 0.5 \\
Total Timesteps & 1,500,000 \\
\hline
\end{tabular}
\end{table}

\subsection{Baseline Methods}

We compare TA-PPO against three baseline routing methods:

\begin{enumerate}
\item \textbf{Dijkstra}: Shortest-path routing based on propagation delay only. Updates routes every 10 timesteps.

\item \textbf{OSPF}: Link-state routing considering both delay and current queue length. Uses cost function: $C = \tau + 2Q$.

\item \textbf{DRL without Prediction}: Same PPO agent but without predicted features in state space (ablation study).
\end{enumerate}

\subsection{Evaluation Metrics}

We evaluate using the following metrics:

\begin{itemize}
\item \textbf{Packet Delivery Ratio (PDR)}: $\frac{\text{Delivered}}{\text{Delivered} + \text{Dropped}} \times 100\%$
\item \textbf{Average Latency}: Mean end-to-end delay for delivered packets
\item \textbf{P95 Latency}: 95th percentile latency
\item \textbf{Average Hop Count}: Mean path length
\item \textbf{Link Utilization}: Mean bandwidth usage across all ISLs
\item \textbf{QoS Satisfaction Rate}: Percentage of packets meeting latency targets
\item \textbf{Drop Reasons}: Breakdown of why packets were dropped
\end{itemize}

All metrics are computed per traffic slice and aggregated globally.

\subsection{Experimental Procedure}

Each experiment consists of:

\begin{enumerate}
\item \textbf{Prediction Model Training}: Train LSTM and GRU on collected data (2.1 minutes for 248 satellites)
\item \textbf{PPO Agent Training}: Train TA-PPO with predictions enabled (1.5M timesteps, ~12 hours)
\item \textbf{Evaluation}: Run each routing method for 5 episodes of 800 timesteps
\item \textbf{Statistical Analysis}: Compute mean and standard deviation across episodes
\end{enumerate}

All experiments were conducted on a machine with:
\begin{itemize}
\item CPU: Intel Core i7-11800H @ 2.3 GHz
\item RAM: 16 GB
\item GPU: Not used (CPU-only for realistic satellite constraints)
\item OS: Windows 10
\item Python: 3.12
\end{itemize}

\subsection{Implementation Details}

Key implementation considerations:

\begin{itemize}
\item \textbf{Prediction Warmup}: First 50 timesteps use baseline routing to collect initial data
\item \textbf{Prediction Caching}: Predictions are cached for 10 timesteps to reduce computation
\item \textbf{Normalization}: All features are normalized to $[0, 1]$ range
\item \textbf{Experience Replay}: PPO uses a buffer of 8192 samples for stable updates
\item \textbf{Early Stopping}: Prediction models use patience of 10 epochs
\end{itemize}

The complete codebase is approximately 5000 lines of Python across multiple modules:

\begin{itemize}
\item \texttt{src/env/topology.py}: Network topology management (500 lines)
\item \texttt{src/routing/drl\_env.py}: PPO environment (800 lines)
\item \texttt{src/prediction/}: Prediction models and data collection (900 lines)
\item \texttt{src/metrics/}: Evaluation and visualization (600 lines)
\item \texttt{src/traffic/}: Traffic generation (400 lines)
\end{itemize}

